\documentclass[11pt]{amsart}
\usepackage{amsmath,amssymb,amsthm}
\usepackage{hyperref}

% Notation
\newcommand{\R}{\mathbb{R}}
\newcommand{\St}{V_k(\R^n)}
\newcommand{\Gr}{\mathrm{Gr}(k,n)}
\DeclareMathOperator{\sym}{sym}
\DeclareMathOperator{\Skew}{skew}
\DeclareMathOperator{\tr}{tr}

% Theorem environments (one shared counter, numbered by section)
\theoremstyle{plain}
\newtheorem{theorem}{Theorem}[section]
\newtheorem{proposition}[theorem]{Proposition}
\newtheorem{lemma}[theorem]{Lemma}
\newtheorem{corollary}[theorem]{Corollary}
\theoremstyle{definition}
\newtheorem{definition}[theorem]{Definition}
\newtheorem{example}[theorem]{Example}
\theoremstyle{remark}
\newtheorem{remark}[theorem]{Remark}

\title{Stiefel and Grassmann Manifolds}
\author{Yuvan Marimuthu}
\date{\today}

\begin{document}

\begin{abstract}
% Write last.
\end{abstract}

\maketitle

\section{Introduction}\label{sec:intro}
% Write last.

\section{Preliminaries}\label{sec:prelim}
Standard references are \cite{lee} and \cite{segal}.

\section{The Stiefel manifold}\label{sec:stiefel}

\section{The Grassmann manifold}\label{sec:grassmann}

\section{From Stiefel to Grassmann}\label{sec:quotient}

\section{Worked example: \texorpdfstring{$\mathrm{Gr}(2,4)$}{Gr(2,4)}}\label{sec:example}

\section{Further directions}\label{sec:further}

\bibliographystyle{amsplain}
\bibliography{refs}

\end{document}
